\section{Determinant-assembly column}
\label{sec:det-assembly}

The four arithmetic columns produce four factors.  To obtain an element of
one determinant line one also needs an orientation: the signs, units and
trivialization conventions under which the factors are multiplied.  The
determinant-assembly column holds this orientation as an independent
input.

\begin{definition}[Real determinant-assembly map]\label{def:apparatus:det-assembly-map}
Let
\[
  \lambda_{\mathrm{loc}}(E)
  =
  \rho_{\mathrm{ht/reg}}^{\R}
  \otimes
  \rho_{\mathrm{an/period}}^{\R}
  \otimes
  \rho_{\mathrm{finite}}^{\R}
  \otimes
  \bigotimes_{p\in S(E)}\rho_p^{\R}
\]
be the tensor of the four arithmetic columns over a finite support set
$S(E)$ of primes, with the local factors at primes outside $S(E)$
trivialized as units.  Let $o_{\det}\in C_{\det}(E)$ be an orientation and
trivialization datum, supplied independently of the four columns, and let
$\rho_{\det}^{\R}(o_{\det})\in\R^\times$ be the orientation scalar it
determines (the map on the orientation slot from
\Cref{def:apparatus:bk-component-map}).  The real determinant assembly is
\[
  \mathcal A^{\R}(E)
  =
  \lambda_{\mathrm{loc}}(E)\cdot\rho_{\det}^{\R}(o_{\det})
  \in\Delta_f(M_E)_{\R}.
\]
We call $\mathcal A^{\R}(E)$ the \emph{final assembly}.  It already
contains the four arithmetic factors, and so must not be used again as a
factor alongside them; see \Cref{def:apparatus:five-column-tensor}.
\end{definition}

The adequacy question for this column is whether the four arithmetic
columns, without the orientation, can explain the assembled target.  In
the coordinate model they cannot, and the shortfall is exactly one
direction.

\begin{theorem}[Cascade-completion Schur signature]\label{thm:apparatus:cascade-completion-signature}
Use the tangent coordinates
\[
  x=(x_{\mathrm{ht}},x_{\mathrm{an}},x_{\mathrm{fin}},x_p,x_{\det})\in\R^5 .
\]
Let the four arithmetic columns supply the source row
$\ell=(1,1,1,1,0)$ and let the assembled target have row
$d=(1,1,1,1,1)$.  With $C=I_5$,
\[
  \KDD=dd^{\mathsf T}=5,\qquad
  \KLL=\ell\ell^{\mathsf T}=4,\qquad
  \KDL=d\ell^{\mathsf T}=4,
  \qquad
  \XiC=5-4\cdot4^{-1}\cdot4=1.
\]
The residual is the orientation direction $e_{\det}=(0,0,0,0,1)$.%
\footnote{Lean proves these exact integer values.  The statement concerns
the orientation direction only; it does not assert adequacy of each of the
four arithmetic coordinates.}
\end{theorem}

\begin{proof}
The three blocks are the displayed sums of squares and products of
entries, and the one-dimensional source block $4$ has inverse $4^{-1}$.
Hence $\XiC=5-4=1$.  Since $d=\ell+e_{\det}$ with $\ell\cdot e_{\det}=0$,
projecting $d$ onto $\ell$ leaves $e_{\det}$, of squared length $1$.
\end{proof}

\begin{obligation}[Integral determinant assembly]\label{obl:apparatus:rho-det-assembly}
The full determinant comparison requires, for each orientation datum
$o_{\det}$, a transport
\[
  \beta_{\det\text{-}\mathrm{assembly}}(\,\cdot\,;o_{\det}):
  \Delta_{\mathrm{ht/reg}}
  \otimes
  \Delta_{\mathrm{an/period}}
  \otimes
  \Delta_{\mathrm{finite}}
  \otimes
  \bigotimes_{p\in S(E)}\Delta_p
  \longrightarrow
  \Delta_f(M_E),
\]
together with compatible unit trivializations of the local lines at primes
outside $S(E)$.  It must fix the Knudsen--Mumford determinant
signs~\citep{KnudsenMumford1976}, the local--global compatibility of
Selmer complexes, the convention relating the zeta element to the
motivic trivialization, and compatibility with the four column
transports.  The orientation datum is a parameter of the transport,
distinct from the assembled output.  This transport is external
arithmetic input.  The residual computation above shows only that, in the
coordinate model, the four arithmetic rows do not account for the
orientation direction; it does not show that the orientation varies
independently for actual curves.
\end{obligation}
